\begin{titlepage}
\thispagestyle{empty}
\begin{tikzpicture}[remember picture,overlay]
  \fill[Navy] (current page.south west) rectangle (current page.north east);
  \fill[Teal] ([xshift=-2.25in]current page.north east) rectangle (current page.north east);
  \foreach \x in {0,...,5}{\foreach \y in {0,...,8}{
    \pgfmathtruncatemacro{\q}{mod(\x+\y,2)}
    \ifnum\q=0\fill[white,opacity=.05] ([xshift=-2.15in+0.36in*\x,yshift=-0.35in-0.36in*\y]current page.north east) rectangle ++(.36in,-.36in);\fi
  }}
  \draw[Gold,line width=2pt] ([xshift=.62in,yshift=-.62in]current page.north west)--([xshift=.62in,yshift=.62in]current page.south west);
\end{tikzpicture}

\vspace*{1.0in}
{\sffamily\bfseries\fontsize{42}{45}\selectfont\color{white} THINKING\\[0.05em] IN CHESS\par}
\vspace{0.35in}
{\sffamily\fontsize{18}{23}\selectfont\color{PaleGold}An Intuitive Game-Theoretic Textbook\par}
\vspace{0.25in}
{\sffamily\large\color{white}How to model positions, read incentives, search trees,\\manage risk, and learn from opponents\par}
\vfill
\begin{tcolorbox}[width=.76\textwidth,colback=white!7!Navy,colframe=white!25!Navy,boxrule=.6pt,arc=1mm,left=4mm,right=4mm,top=3mm,bottom=3mm]
{\color{white}\small\sffamily
A chess move is not an isolated object. It is an action inside a living tree of replies. This book teaches you to see that tree without drowning in it.}
\end{tcolorbox}
\vspace{0.55in}
{\sffamily\color{white!75}\small Original educational synthesis prepared from the game-theory framework supplied by the reader\par}
\end{titlepage}

\cleardoublepage
\thispagestyle{empty}
\vspace*{1.0in}
{\Large\sffamily\bfseries\color{Navy}Thinking in Chess\par}
\vspace{0.35in}
\textbf{Edition 1.0}

This is an original textbook. It does not reproduce the wording, examples, exercises, or chapter text of the uploaded game-theory book. Instead, it translates a broad sequence of ideas---game frames, preferences and utility, strategic choice, extensive-form reasoning, uncertainty, beliefs, and incomplete information---into a practical language for chess.

All chess diagrams and examples in this book are pedagogical constructions unless otherwise stated. Piece values and probability estimates are heuristics, not laws. In any legal dispute over a position, the official rules of the event govern.

\vfill
\begin{center}
{\sffamily\small\color{Teal}MODEL \quad REPLY \quad VALUE \quad DECIDE \quad UPDATE}
\end{center}
\cleardoublepage

\chapter*{Preface: The move is not the point}
\addcontentsline{toc}{chapter}{Preface: The move is not the point}

Most chess instruction presents answers: a tactic, a rule, an opening move, a famous plan. Answers are useful, but they decay when the position changes. A way of thinking travels farther.

Game theory begins from a deceptively simple observation: your choice cannot be evaluated in isolation from someone else's choice. Chess makes that observation visible. Every move changes the actions available to your opponent, the value of future positions, and the information both players reveal about themselves. The board is fully visible, yet the decision remains uncertain because neither player can calculate everything and neither knows exactly what the other has noticed.

This book therefore treats chess as a sequence of structured decisions. You will learn to ask five questions:

\begin{enumerate}[label=\textbf{\arabic*.},leftmargin=2.2em]
\item \textbf{What is the state?} Record material, safety, activity, structure, time, and whose move it is.
\item \textbf{What am I trying to maximize?} The answer is not always ``material.'' It may be winning chances, drawing security, simplicity, time, or tournament score.
\item \textbf{What can both players do?} Generate a small set of candidates, then search for the opponent's strongest reply.
\item \textbf{How do the branches compare?} Use worst-case reasoning for forcing lines and expected-value reasoning for practical uncertainty.
\item \textbf{What did the last move teach me?} Update your model of the position and of the opponent.
\end{enumerate}

The result is a compact loop:

\begin{center}
\begin{tikzpicture}[node distance=10mm and 7mm,>=Latex]
\tikzset{flowstep/.style={draw=Teal,rounded corners=2mm,fill=PaleTeal,minimum width=22mm,minimum height=9mm,font=\sffamily\bfseries\small,text=Navy}}
\node[flowstep] (m) {MODEL};
\node[flowstep,right=of m] (r) {REPLY};
\node[flowstep,right=of r] (v) {VALUE};
\node[flowstep,right=of v] (d) {DECIDE};
\node[flowstep,below=of v] (u) {UPDATE};
\draw[->,thick,Teal] (m)--(r);
\draw[->,thick,Teal] (r)--(v);
\draw[->,thick,Teal] (v)--(d);
\draw[->,thick,Teal] (d) |- (u);
\draw[->,thick,Teal] (u) -| (m);
\end{tikzpicture}
\end{center}

You do not need advanced mathematics. Formal boxes make the logic precise, but every formal idea is paired with an intuition, a chess example, and a memory hook. The aim is not to turn you into an engine. It is to help you allocate human attention well.

\section*{Who this book is for}
The core reader knows how the pieces move and has encountered basic game-theory words such as strategy, payoff, and Nash equilibrium. A rules refresher is included. Beginners can read straight through; experienced players can use the chapters on minimax, uncertainty, opponent modeling, time management, and review as a separate course.

\section*{How to use it}
Read one chapter at a time. Before checking a worked example, cover the continuation and name the opponent's best reply. Write answers to the exercises. During your next games, choose only one chapter habit to practice. A habit applied at the board is worth more than ten ideas recognized on the page.

The appendices contain a notation guide, printable decision sheets, complete exercise solutions, and a glossary. Chapter 26 turns the book into a twelve-week study plan.

\oneminute{Chess improvement is not mainly the accumulation of more moves. It is the improvement of the questions you ask before choosing one.}

\chapter*{A map from game theory to chess}
\addcontentsline{toc}{chapter}{A map from game theory to chess}

\begin{longtable}{@{}p{.25\textwidth}p{.28\textwidth}p{.34\textwidth}@{}}
\toprule
\textbf{Game-theory lens} & \textbf{Chess translation} & \textbf{Practical question}\\
\midrule
Game frame & Board, side to move, legal moves, clocks, rules & What facts define the decision?\\
Preferences and utility & Win/draw/loss plus practical goals & What result and what kind of position do I need?\\
Strategic form & Complete plans and opening repertoires & What will I do against each broad response?\\
Extensive form & The move tree & What reply follows each candidate?\\
Best response & Strongest counterplay & What would I hate to see after my move?\\
Dominance & One candidate is no worse in relevant branches & Can I safely eliminate a move?\\
Nash equilibrium & Mutually optimal strategies & What survives best play from both sides?\\
Backward induction & Solve from the leaf positions backward & Which forcing ending settles the line?\\
Expected utility & Value weighted by practical likelihood & Which move scores best against a fallible human?\\
Mixed strategy & Varying preparation across repeated games & Am I becoming predictable?\\
Beliefs and Bayesian updating & Model of opponent's style and knowledge & What did the last move reveal?\\
Incomplete information & Hidden preparation, intentions, fatigue & Which unknowns should affect my choice?\\
Repeated games & Matches, tournaments, reputation & How does today's decision change tomorrow's incentives?\\
Bounded rationality & Limited time, memory, and calculation & Where should I spend attention?\\
\bottomrule
\end{longtable}

\begin{warningbox}[title={A crucial distinction}]
Chess is a perfect-information game because the board and legal history are observable. Human chess is still a decision problem under uncertainty because calculation, memory, intentions, and future errors are not observable. Do not confuse hidden minds with hidden pieces.
\end{warningbox}

\cleardoublepage
